%%%%%%%%%%%%%%%%%%%%%%%%%%%%%%%%%%%%%%%%%%%%%%%%%%%%%%%%%%%%%%%%%%%%%%%%%%
%
% Copyright (c) 2008-2011, Nokia Corporation and/or its subsidiary(-ies).
% All rights reserved.
%
% This work, unless otherwise expressly stated, is licensed under a
% Creative Commons Attribution-ShareAlike 2.5.
%
% The full license document is available from
% http://creativecommons.org/licenses/by-sa/2.5/legalcode .
%
%%%%%%%%%%%%%%%%%%%%%%%%%%%%%%%%%%%%%%%%%%%%%%%%%%%%%%%%%%%%%%%%%%%%%%%%%%

\documentclass[t]{beamer}
%%%%%%%%%%%%%%%%%%%%%%%%%%%%%%%%%%%%%%%%%%%%%%%%%%%%%%%%%%%%%%%%%%%%%%%%%%
%
% This work, unless otherwise expressly stated, is licensed under a
% Creative Commons Attribution-ShareAlike 2.5.
%
% The full license document is available from
% http://creativecommons.org/licenses/by-sa/2.5/legalcode .
%
%%%%%%%%%%%%%%%%%%%%%%%%%%%%%%%%%%%%%%%%%%%%%%%%%%%%%%%%%%%%%%%%%%%%%%%%%%

\usepackage{ifthen}

\newcommand{\coursetitle}{\title}
\newcommand{\producer}{\subtitle}
\newcommand{\material}{\institute}
\newcommand{\customer}{\author}
\newcommand{\partner}{\titlegraphic}

\newboolean{eval-version}
\newboolean{with-notes}
\newboolean{show-slide-id}
\newboolean{show-page-numbers}


%======================================================================
%                          Configuration of Title Page
%======================================================================


% General course title
\coursetitle{Qt Essentials}
\producer{\small{Training Course}}

% Presented for Customer or Presented by Trainer
\customer{Visit us at \url{http://qt-project.org}}

% Indicator for material version
\material{
  Produced by the Qt Training Partners \\
  \medskip
  \tiny{\textit{Material based on Qt 5.1, created on \today\\}}
}

% Partner company logo and name
% \partner{\includegraphics[height=2cm]{images/company} \\ The Partner}

%======================================================================
%                          Configuration of ifdefs
%======================================================================

%% Include ``Evaluation version'' on the slides?
\setboolean{eval-version}{false}

%% Whether to include instructors notes
\setboolean{with-notes}{false}

%% Should we show page numbers on in the footer?
\setboolean{show-page-numbers}{true}

%% Should we show slide-id's in the title?
\setboolean{show-slide-id}{false}

%% What is the screen resolution
\newcommand{\screenwidth}{1400}

%%%%%%%%%%%%%%%%%%%%%%%%%%%%%%%%%%%%%%%%%%%%%%%%%%%%%%%%%%%%%%%%%%%%%%%%%%
%
% Copyright (c) 2008-2011, Nokia Corporation and/or its subsidiary(-ies).
% All rights reserved.
%
% This work, unless otherwise expressly stated, is licensed under a
% Creative Commons Attribution-ShareAlike 2.5.
%
% The full license document is available from
% http://creativecommons.org/licenses/by-sa/2.5/legalcode .
%
%%%%%%%%%%%%%%%%%%%%%%%%%%%%%%%%%%%%%%%%%%%%%%%%%%%%%%%%%%%%%%%%%%%%%%%%%%

\usepackage{ifthen}
\usepackage{alltt}
\usepackage{comment}
\usepackage{multicol}
\usepackage{multirow}
\usepackage{hhline}
\usepackage{bold-extra}
\usepackage{version}
\usepackage{graphicx}
\usepackage{pgfpages}
\usepackage{picins}
\usepackage{makeidx}
\usepackage{multind}
\usepackage{xspace}
\usepackage{mathptmx}
\usepackage{pdfsync}
\usepackage{listings}
%\usepackage{multimedia}

\ifthenelse{\boolean{cc-license}}
{
  \usetheme{QtTraining}
} % else
{
  \usetheme{QtNokia}
}

% \definecolor{code_bg}{RGB}{255, 255, 204}

\lstdefinelanguage{qmake}%
  { morekeywords={SOURCES,HEADERS,RESOURCES,FORMS},
    morestring=[b]",
    sensitive=true,
    comment=[l]\#,
 }

\ifthenelse{\boolean{cc-license}}
{
  \lstdefinelanguage[Qt]{C++}[]{C++}%
{%
  morekeywords=[2]{%
    Q_OBJECT,Q_PROPERTY,%
    SIGNAL,SLOT,%
    signals,slots,connect,disconnect,%
    emit,tr,qobject_cast,foreach,%
    qscriptvalue_cast%
  },%
  %
  % Generated from Qt 4.7 include dir:
  % $ find . -name 'Q*' -printf '%f,%%\n'
  morekeywords=[3]{%
    QtUiTools,%
    QtUiTools,%
    QUiLoader,%
    QtNetwork,%
    QNetworkInterface,%
    QHostAddress,%
    QNetworkProxy,%
    QtNetwork,%
    QSsl,%
    QAuthenticator,%
    QSslCipher,%
    QNetworkAccessManager,%
    QAbstractSocket,%
    QNetworkRequest,%
    QLocalSocket,%
    QLocalServer,%
    QHttp,%
    QUdpSocket,%
    QTcpServer,%
    QTcpSocket,%
    QHttpResponseHeader,%
    QNetworkCookie,%
    QNetworkConfigurationManager,%
    QIPv6Address,%
    QSslCertificate,%
    QNetworkProxyFactory,%
    QUrlInfo,%
    QNetworkProxyQuery,%
    QNetworkAddressEntry,%
    QNetworkConfiguration,%
    QHostInfo,%
    QFtp,%
    QSslKey,%
    QSslError,%
    QNetworkCookieJar,%
    QNetworkReply,%
    QHttpHeader,%
    Q_IPV6ADDR,%
    QNetworkCacheMetaData,%
    QSslConfiguration,%
    QNetworkSession,%
    QNetworkDiskCache,%
    QSslSocket,%
    QHttpRequestHeader,%
    QAbstractNetworkCache,%
    Qt,%
    QtCore,%
    QMetaMethod,%
    QTextIStream,%
    QListIterator,%
    QBitArray,%
    QCoreApplication,%
    Q_INT32,%
    QForeachContainerBase,%
    QTextStream,%
    QFile,%
    QMetaObjectExtraData,%
    QFileInfoListIterator,%
    QThreadStorage,%
    QGlobalStatic,%
    QSysInfo,%
    QByteArray,%
    QStdWString,%
    QDebug,%
    QMutableFutureIterator,%
    QSettings,%
    Qt,%
    QAtomicPointer,%
    QProcess,%
    QLatin1Literal,%
    QThreadPool,%
    QChildEvent,%
    QReadLocker,%
    Q_LLONG,%
    QObjectCleanupHandler,%
    QtCore,%
    QHashIterator,%
    QStringRef,%
    QLocale,%
    Q_LONG,%
    QList,%
    QDir,%
    QResource,%
    QAnimationGroup,%
    QHashData,%
    QInternal,%
    QXmlStreamReader,%
    QMutableSetIterator,%
    QFuture,%
    QSetIterator,%
    QAbstractFileEngine,%
    QAbstractTableModel,%
    QRunnable,%
    QMapNode,%
    QtCleanUpFunction,%
    QStringList,%
    QEventTransition,%
    QTextStreamManipulator,%
    QPointer,%
    QContiguousCacheData,%
    Q_INT8,%
    QVector,%
    QXmlStreamEntityDeclarations,%
    QSharedDataPointer,%
    QtPlugin,%
    QTextBoundaryFinder,%
    QDateTime,%
    QMutableMapIterator,%
    QWeakPointer,%
    QPauseAnimation,%
    QVariantMap,%
    QExplicitlySharedDataPointer,%
    QXmlStreamWriter,%
    QStateMachine,%
    QDynamicPropertyChangeEvent,%
    QXmlStreamAttributes,%
    QPluginLoader,%
    QSize,%
    QPointF,%
    QBool,%
    QMetaObjectAccessor,%
    QLinkedListData,%
    Q_INT16,%
    QFinalState,%
    QSet,%
    QTextOStream,%
    QtConcurrentFilter,%
    QTextCodec,%
    QAbstractFileEngineIterator,%
    QSharedMemory,%
    QXmlStreamNamespaceDeclaration,%
    QVariantHash,%
    QHashDummyValue,%
    QReadWriteLock,%
    QObject,%
    QLinkedListIterator,%
    QSocketNotifier,%
    QConcatenable,%
    QMutableLinkedListIterator,%
    QVariant,%
    QTypeInfo,%
    QMargins,%
    QSystemSemaphore,%
    QIncompatibleFlag,%
    QMutableVectorIterator,%
    QtDebug,%
    QSequentialAnimationGroup,%
    QPersistentModelIndex,%
    Q_UINT32,%
    Q_ULLONG,%
    QScopedPointerDeleter,%
    QFutureInterface,%
    QMutableListIterator,%
    QMimeData,%
    QWaitCondition,%
    QBasicAtomicPointer,%
    QFileInfo,%
    QFileInfoList,%
    QTimeLine,%
    QMetaProperty,%
    QGenericArgument,%
    QTextDecoder,%
    QFactoryInterface,%
    QXmlStreamEntityResolver,%
    QVariantList,%
    QString,%
    QEasingCurve,%
    QMap,%
    QNoImplicitBoolCast,%
    QTemporaryFile,%
    QLinkedList,%
    QStringBuilder,%
    QTextEncoder,%
    QMetaObject,%
    QThreadStorageData,%
    QIODevice,%
    QAbstractFileEngineHandler,%
    QtEndian,%
    QVariantAnimation,%
    QScopedPointer,%
    QMetaClassInfo,%
    QFileSystemWatcher,%
    QLineF,%
    QtConfig,%
    QContiguousCacheTypedData,%
    QProcessEnvironment,%
    QForeachContainer,%
    QTS,%
    QVectorIterator,%
    QRectF,%
    QAtomicInt,%
    QtGlobal,%
    QVarLengthArray,%
    QByteRef,%
    QtConcurrentRun,%
    QCOORD,%
    QLatin1Char,%
    QStringListIterator,%
    QReturnArgument,%
    QDataStream,%
    QEventLoop,%
    QModelIndex,%
    QConstString,%
    QMetaTypeId,%
    QLibrary,%
    QParallelAnimationGroup,%
    QSharedPointer,%
    QVariantComparisonHelper,%
    QObjectUserData,%
    QTimer,%
    QCryptographicHash,%
    QThread,%
    QStringMatcher,%
    QXmlStreamAttribute,%
    QWriteLocker,%
    QHistoryState,%
    QQueue,%
    QtMsgHandler,%
    QElapsedTimer,%
    QSignalMapper,%
    QFSFileEngine,%
    QHashNode,%
    QUrl,%
    QFutureWatcherBase,%
    Q_ULONG,%
    QSizeF,%
    QAbstractConcatenable,%
    QByteArrayMatcher,%
    QFutureInterfaceBase,%
    QAbstractItemModel,%
    QTranslator,%
    QtPluginInstanceFunction,%
    QGlobalStaticDeleter,%
    QtConcurrentMap,%
    QNoDebug,%
    Q_UINT64,%
    QAbstractAnimation,%
    QObjectList,%
    QScopedPointerPodDeleter,%
    QMutexLocker,%
    QGenericReturnArgument,%
    QRect,%
    QPropertyAnimation,%
    QRegExp,%
    QMapData,%
    QXmlStreamNotationDeclaration,%
    QMultiMap,%
    QFutureSynchronizer,%
    QDirIterator,%
    QTextCodecPlugin,%
    QXmlStreamStringRef,%
    QMetaTypeId2,%
    QPair,%
    QHashDummyNode,%
    QLatin1String,%
    QModelIndexList,%
    QObjectData,%
    QXmlStreamNotationDeclarations,%
    QBasicAtomicInt,%
    QMutableHashIterator,%
    QMapPayloadNode,%
    QPoint,%
    QDate,%
    QSemaphore,%
    QIntegerForSize,%
    QLibraryInfo,%
    QBitRef,%
    QHash,%
    QMutex,%
    QArgument,%
    QMetaType,%
    QCache,%
    QtContainerFwd,%
    QTextCodecFactoryInterface,%
    QMutableStringListIterator,%
    QFutureWatcher,%
    QLinkedListNode,%
    Q_INT64,%
    QChar,%
    QFlags,%
    QUuid,%
    QLine,%
    QBasicTimer,%
    QStack,%
    QEvent,%
    QFutureIterator,%
    QScopedPointerArrayDeleter,%
    QVectorData,%
    QFlag,%
    Q_UINT16,%
    QSignalTransition,%
    QAbstractEventDispatcher,%
    QMapIterator,%
    QListData,%
    QMultiHash,%
    QTimerEvent,%
    QAbstractListModel,%
    Q_UINT8,%
    QState,%
    QScopedArrayPointer,%
    Q_PID,%
    QAbstractState,%
    QXmlStreamNamespaceDeclarations,%
    QTime,%
    QtAlgorithms,%
    QBuffer,%
    QCharRef,%
    QSharedData,%
    QVectorTypedData,%
    QSystemLocale,%
    QTextStreamFunction,%
    QAbstractTransition,%
    QCustomEvent,%
    QMetaEnum,%
    QContiguousCache,%
    QXmlStreamEntityDeclaration,%
    QtDesigner,%
    QDesignerWidgetDataBaseItemInterface,%
    QDesignerDynamicPropertySheetExtension,%
    QDesignerLanguageExtension,%
    QDesignerMetaDataBaseItemInterface,%
    QDesignerCustomWidgetCollectionInterface,%
    QtDesigner,%
    QDesignerIntegrationInterface,%
    QDesignerBrushManagerInterface,%
    QDesignerWidgetFactoryInterface,%
    QDesignerMemberSheetExtension,%
    QDesignerWidgetDataBaseInterface,%
    QDesignerDnDItemInterface,%
    QDesignerResourceBrowserInterface,%
    QDesignerPropertyEditorInterface,%
    QAbstractFormBuilder,%
    QDesignerFormWindowToolInterface,%
    QDesignerExtraInfoExtension,%
    QExtensionFactory,%
    QDesignerWidgetBoxInterface,%
    QDesignerActionEditorInterface,%
    QAbstractExtensionFactory,%
    QDesignerFormWindowManagerInterface,%
    QExtensionManager,%
    QDesignerFormWindowInterface,%
    QDesignerFormEditorInterface,%
    QDesignerFormEditorPluginInterface,%
    QFormBuilder,%
    QDesignerMetaDataBaseInterface,%
    QDesignerPromotionInterface,%
    QDesignerObjectInspectorInterface,%
    QDesignerExportWidget,%
    QDesignerLayoutDecorationExtension,%
    QDesignerPropertySheetExtension,%
    QDesignerFormWindowCursorInterface,%
    QAbstractExtensionManager,%
    QDesignerContainerExtension,%
    QDesignerTaskMenuExtension,%
    QDesignerComponents,%
    QDesignerCustomWidgetInterface,%
    QDesignerIconCacheInterface,%
    QtMultimedia,%
    QAudio,%
    QAbstractVideoSurface,%
    QAudioInput,%
    QAbstractVideoBuffer,%
    QtMultimedia,%
    QAbstractAudioInput,%
    QAudioEngineFactoryInterface,%
    QVideoSurfaceFormat,%
    QAbstractAudioDeviceInfo,%
    QAudioDeviceInfo,%
    QAudioFormat,%
    QAudioOutput,%
    QVideoFrame,%
    QAbstractAudioOutput,%
    QAudioEnginePlugin,%
    QAxFactory,%
    QAxSelect,%
    QAxObject,%
    QAxScriptManager,%
    QAxAggregated,%
    QAxBindable,%
    QAxScript,%
    QAxBase,%
    QAxScriptEngine,%
    QAxClass,%
    QAxWidget,%
    QtScript,%
    QScriptExtensionInterface,%
    QScriptContextInfo,%
    QScriptValueList,%
    QScriptProgram,%
    QtScript,%
    QScriptEngineAgent,%
    QScriptValue,%
    QScriptEngine,%
    QScriptable,%
    QScriptClassPropertyIterator,%
    QScriptContext,%
    QScriptExtensionPlugin,%
    QScriptSyntaxCheckResult,%
    QScriptValueIterator,%
    QScriptString,%
    QScriptContextInfoList,%
    QScriptClass,%
    QtMeeGoGraphicsSystemHelper,%
    QMeeGoLivePixmap,%
    QMeeGoOverlayWidget,%
    QMeeGoRuntime,%
    QtMeeGoGraphicsSystemHelper,%
    QMeeGoGraphicsSystemHelper,%
    QMeeGoLiveImage,%
    QtScriptTools,%
    QtScriptTools,%
    QScriptEngineDebugger,%
    QtSvg,%
    QSvgGenerator,%
    QSvgWidget,%
    QGraphicsSvgItem,%
    QtSvg,%
    QSvgRenderer,%
    QtSql,%
    QPSQLDriver,%
    QODBCResult,%
    QSqlDriver,%
    QSqlQuery,%
    QSqlResult,%
    QSqlTableModel,%
    QSqlRelationalTableModel,%
    QSQLiteDriver,%
    QIBaseDriver,%
    QSqlDriverCreatorBase,%
    QSQLiteResult,%
    QDB2Driver,%
    QMYSQLResult,%
    QtSql,%
    QSqlDatabase,%
    QTDSResult,%
    QSqlRelation,%
    QDB2Result,%
    QSqlIndex,%
    QIBaseResult,%
    QSqlField,%
    QSqlRecord,%
    QODBCDriver,%
    QSqlDriverFactoryInterface,%
    QPSQLResult,%
    QMYSQLDriver,%
    QSQLite2Result,%
    QSqlRelationalDelegate,%
    QOCIDriver,%
    QSQLite2Driver,%
    QOCIResult,%
    QSqlQueryModel,%
    QSqlDriverPlugin,%
    QTDSDriver,%
    QSqlDriverCreator,%
    QSqlError,%
    QtDBus,%
    QDBusReply,%
    QDBusServiceWatcher,%
    QDBusConnection,%
    QDBusError,%
    QDBusSignature,%
    QDBusMessage,%
    QDBusAbstractInterfaceBase,%
    QDBusContext,%
    QDBusServer,%
    QDBusConnectionInterface,%
    QDBusInterface,%
    QDBusPendingCallWatcher,%
    QDBusAbstractAdaptor,%
    QDBusPendingReplyData,%
    QDBusArgument,%
    QtDBus,%
    QDBusMetaType,%
    QDBusPendingCall,%
    QDBusVariant,%
    QDBusObjectPath,%
    QDBusAbstractInterface,%
    QDBusPendingReply,%
    QtGui,%
    QDragResponseEvent,%
    QWheelEvent,%
    QPrintDialog,%
    QTreeWidget,%
    QItemSelectionRange,%
    QGraphicsGridLayout,%
    QStandardItem,%
    QCalendarWidget,%
    QStyleOptionFocusRect,%
    QFormLayout,%
    QLinearGradient,%
    QItemSelectionModel,%
    QGraphicsColorizeEffect,%
    QWSServer,%
    QDialog,%
    QKeySequence,%
    QTextTable,%
    QWSKeyboardHandler,%
    QSyntaxHighlighter,%
    QRubberBand,%
    QAction,%
    QWSSoundServerSocket,%
    QWSInputMethod,%
    QItemEditorCreator,%
    QApplication,%
    QAuthDevice,%
    QTextBlockFormat,%
    QGraphicsEllipseItem,%
    QStyleOptionTabWidgetFrameV2,%
    QWSCursor,%
    QScrollBar,%
    QWSEvent,%
    QMouseDriverFactory,%
    QKbdDriverPlugin,%
    QUndoView,%
    QGradient,%
    QGraphicsRectItem,%
    QS60StubMAknTouchPaneObserver,%
    QGraphicsSceneContextMenuEvent,%
    QGroupBox,%
    QStyleOptionTabBarBaseV2,%
    QCDEStyle,%
    QPicture,%
    QScreenDriverFactory,%
    QColormap,%
    QStylePainter,%
    QIconDragEvent,%
    QWSQnxKeyboardHandler,%
    QWSScreenSaver,%
    QTableView,%
    QPrinterInfo,%
    QGraphicsPolygonItem,%
    QMotifStyle,%
    QTabBar,%
    QAccessibleInterfaceEx,%
    QFrame,%
    QTextLine,%
    QAccessibleSimpleEditableTextInterface,%
    QPolygonF,%
    QGraphicsOpacityEffect,%
    QMovie,%
    QDesktopServices,%
    QGraphicsLineItem,%
    QAccessibleFactoryInterface,%
    QPaintEngine,%
    QPixmapCache,%
    QCommonStyle,%
    QMatrix3x4,%
    QWhatsThisClickedEvent,%
    QInputEvent,%
    QGraphicsRotation,%
    QGraphicsSceneHelpEvent,%
    QSystemTrayIcon,%
    QStyleOptionViewItemV3,%
    QGraphicsTransform,%
    QLinuxFb_Shared,%
    QMatrix2x4,%
    QGraphicsAnchor,%
    QProgressBar,%
    QPrintEngine,%
    QSizePolicy,%
    QMenubarUpdatedEvent,%
    QMenuItem,%
    QContextMenuEvent,%
    QAccessibleActionInterface,%
    QPushButton,%
    QStyledItemDelegate,%
    QWSCursorMap,%
    QTapAndHoldGesture,%
    QDirectPainter,%
    QWidgetSet,%
    QItemSelection,%
    QS60MainApplicationBase,%
    QTableWidget,%
    QAbstractScrollArea,%
    QPaintEvent,%
    QAccessibleBridge,%
    QTextTableCell,%
    QToolTip,%
    QMimeSource,%
    QDecorationDefault,%
    QDoubleSpinBox,%
    QVector3D,%
    QGraphicsView,%
    QDragMoveEvent,%
    QVFbMouseHandler,%
    QStyleOptionDockWidget,%
    QGraphicsDropShadowEffect,%
    QTextBrowser,%
    QColorDialog,%
    QCompleter,%
    QDecoration,%
    QWindowsStyle,%
    QMenuBar,%
    QAccessibleWidgetEx,%
    QPinchGesture,%
    QLayoutIterator,%
    QWidget,%
    QImageReader,%
    QToolBar,%
    QMacMime,%
    QPlainTextDocumentLayout,%
    QImageWriter,%
    QDial,%
    QGtkStyle,%
    QStyleOptionMenuItem,%
    QWindowsMobileStyle,%
    QStyleOptionQ3DockWindow,%
    QInputContext,%
    QVFbKeyboardHandler,%
    QTextBlockGroup,%
    QColor,%
    QStyleFactoryInterface,%
    QTreeView,%
    QGraphicsPathItem,%
    QTreeWidgetItem,%
    QPrintPreviewWidget,%
    QInputContextPlugin,%
    QToolBarChangeEvent,%
    QPixmap,%
    QStyleOptionButton,%
    QBoxLayout,%
    QWSServerSocket,%
    QActionGroup,%
    QWSPropertyManager,%
    QStyleOptionProgressBarV2,%
    QDateTimeEdit,%
    QX11EmbedContainer,%
    QTextCharFormat,%
    QClipboardEvent,%
    QKbdDriverFactory,%
    QScrollArea,%
    QDragEnterEvent,%
    QStyleOptionQ3ListView,%
    QTextBlockUserData,%
    QTextImageFormat,%
    QStyleOptionComboBox,%
    QVector4D,%
    QLayoutItem,%
    QStyleOptionTabBarBase,%
    QMacNativeWidget,%
    QValidator,%
    QInputDialog,%
    QStyleOptionGraphicsItem,%
    QStandardItemModel,%
    QTextEdit,%
    QWSClient,%
    QStyleOptionSlider,%
    QPrinter,%
    QTextDocumentFragment,%
    QStyleOptionToolButton,%
    QWSSoundClient,%
    QTableWidgetItem,%
    QMatrix3x3,%
    QGraphicsObject,%
    QCommandLinkButton,%
    QS60MainApplication,%
    QStyleOptionComplex,%
    QWindowStateChangeEvent,%
    QAbstractSpinBox,%
    QFileSystemModel,%
    QGraphicsSceneMoveEvent,%
    QTreeWidgetItemIterator,%
    QScreen,%
    QDoubleValidator,%
    QTextListFormat,%
    QScreenDriverFactoryInterface,%
    QS60StubAknAppUi,%
    QMatrix2x3,%
    QDropEvent,%
    QGraphicsLayout,%
    QGestureRecognizer,%
    QGraphicsLinearLayout,%
    QStyleOptionToolBoxV2,%
    QAccessibleWidget,%
    QTransportAuth,%
    QWidgetItem,%
    QDragLeaveEvent,%
    QClipboard,%
    QHeaderView,%
    QTextDocumentWriter,%
    QPalette,%
    QItemEditorFactory,%
    QDecorationFactory,%
    QMainWindow,%
    QGraphicsSceneMouseEvent,%
    QColumnView,%
    QMatrix,%
    QGraphicsSceneEvent,%
    QMdiArea,%
    QTextObjectInterface,%
    QTextLayout,%
    QIcon,%
    QStyleOptionFrameV3,%
    QFocusEvent,%
    QAbstractGraphicsShapeItem,%
    QPainterPathStroker,%
    QShortcutEvent,%
    QSplitterHandle,%
    QTextTableCellFormat,%
    QTextBlock,%
    QRadioButton,%
    QAccessibleObject,%
    QIconEngineFactoryInterface,%
    QStyleOptionViewItemV2,%
    QGesture,%
    QStyleOptionToolBar,%
    QLCDNumber,%
    QProxyModel,%
    QSplitter,%
    QLineEdit,%
    QDialogButtonBox,%
    QFocusFrame,%
    QStyle,%
    QS60Style,%
    QFontInfo,%
    QWSProtocolItem,%
    QCursor,%
    QScreenDriverPlugin,%
    QAccessibleImageInterface,%
    QGradientStop,%
    QS60MainAppUiBase,%
    QDateEdit,%
    QTabletEvent,%
    QScreenCursor,%
    QTextFrame,%
    QWindowsCEStyle,%
    QWSMouseHandlerFactoryInterface,%
    QLabel,%
    QDecorationWindows,%
    QAbstractItemView,%
    QStyleOptionHeader,%
    QProgressDialog,%
    QPolygon,%
    QAccessibleTableInterface,%
    QStyleOptionTitleBar,%
    QFontMetricsF,%
    QMenu,%
    QGraphicsItemAnimation,%
    QWidgetList,%
    QAbstractPageSetupDialog,%
    QIntValidator,%
    QGraphicsSceneWheelEvent,%
    QTextInlineObject,%
    QStyleOptionTabWidgetFrame,%
    QWorkspace,%
    QPictureFormatInterface,%
    QWSUmKeyboardHandler,%
    QItemDelegate,%
    QAccessibleTextInterface,%
    QPaintEngineState,%
    QGenericMatrix,%
    QHBoxLayout,%
    QBrushData,%
    QRegExpValidator,%
    QX11EmbedWidget,%
    QAbstractFontEngine,%
    QAccessibleObjectEx,%
    QTextFormat,%
    QMatrix4x2,%
    QTextFragment,%
    QBrush,%
    QTransform,%
    QS60MainAppUi,%
    QAbstractSlider,%
    QStyleOptionTabV3,%
    QListWidgetItem,%
    QTapGesture,%
    QUndoGroup,%
    QTextOption,%
    QResizeEvent,%
    QStyleOptionViewItemV4,%
    QImage,%
    QWSMouseHandler,%
    QWSLinuxInputKeyboardHandler,%
    QX11Info,%
    QGraphicsAnchorLayout,%
    QFileOpenEvent,%
    QImageIOHandlerFactoryInterface,%
    QTextCursor,%
    QDecorationPlugin,%
    QVector2D,%
    QDrag,%
    QPainterPath,%
    QItemEditorCreatorBase,%
    QAccessibleEvent,%
    QStatusTipEvent,%
    QStylePlugin,%
    QStyleFactory,%
    QPainter,%
    QAccessibleApplication,%
    QListWidget,%
    QWMatrix,%
    QAccessible2Interface,%
    QTextTableFormat,%
    QTextLength,%
    QAccessibleInterface,%
    QProxyScreen,%
    QUnixPrintWidget,%
    QUndoStack,%
    QTransformedScreen,%
    QS60MainDocumentBase,%
    QSplashScreen,%
    QUpdateLaterEvent,%
    QDecorationFactoryInterface,%
    QAccessibleBridgeFactoryInterface,%
    QStyleHintReturnVariant,%
    QImageIOHandler,%
    QStyleOptionViewItem,%
    QWidgetAction,%
    QStandardItemEditorCreator,%
    QGraphicsLayoutItem,%
    QKeyEventTransition,%
    QAbstractPrintDialog,%
    QStyleOptionDockWidgetV2,%
    QAbstractProxyModel,%
    QS60StubMEikStatusPaneObserver,%
    QSwipeGesture,%
    QGestureEvent,%
    QShortcut,%
    QMacPasteboardMime,%
    QTabWidget,%
    QMouseDriverPlugin,%
    QMacStyle,%
    QAbstractTextDocumentLayout,%
    QWSInternalWindowInfo,%
    QWSDisplay,%
    QTextObject,%
    QTextFrameFormat,%
    QTableWidgetSelectionRange,%
    QFontEngineInfo,%
    QMatrix3x2,%
    QAccessibleValueInterface,%
    QPaintDevice,%
    QMouseEventTransition,%
    QWSKeyboardHandlerFactoryInterface,%
    QSound,%
    QGraphicsItemGroup,%
    QGraphicsSceneHoverEvent,%
    QAbstractUndoItem,%
    QSizeGrip,%
    QRgb,%
    QWidgetData,%
    QVFbHeader,%
    QStyleOptionRubberBand,%
    QAccessiblePlugin,%
    QGraphicsWidget,%
    QIconSet,%
    QStyleOptionFrameV2,%
    QPlainTextEdit,%
    QToolBox,%
    QHideEvent,%
    QCopChannel,%
    QStackedLayout,%
    QProxyScreenCursor,%
    QCleanlooksStyle,%
    QWSManager,%
    QMessageBox,%
    QWSSocket,%
    QTextItem,%
    QStyleOptionToolBox,%
    QWidgetItemV2,%
    QS60StubAknAppUiBase,%
    QStaticText,%
    QQnxMouseHandler,%
    QStyleOptionProgressBar,%
    QPictureIO,%
    QConicalGradient,%
    QPen,%
    QS60MainDocument,%
    QMacCocoaViewContainer,%
    QGridLayout,%
    QIconEnginePlugin,%
    QtGui,%
    QWSPointerCalibrationData,%
    QFont,%
    QSessionManager,%
    QRegion,%
    QGraphicsSimpleTextItem,%
    QStyleOption,%
    QWSWindow,%
    QImageIOPlugin,%
    QMatrix4x4,%
    QGradientStops,%
    QLayout,%
    QWhatsThis,%
    QMoveEvent,%
    QIconEngineFactoryInterfaceV2,%
    QTextDocument,%
    QCloseEvent,%
    QInputContextFactory,%
    QWizard,%
    QWSCalibratedMouseHandler,%
    QQuaternion,%
    QInputContextFactoryInterface,%
    QGraphicsEffect,%
    QPoolEntry,%
    QFileIconProvider,%
    QIconEngineV2,%
    QShowEvent,%
    QAbstractItemDelegate,%
    QQnxScreen,%
    QHelpEvent,%
    QDockWidget,%
    QStyleOptionFrame,%
    QSlider,%
    QWSWindowInfo,%
    QHoverEvent,%
    QWizardPage,%
    QVFbScreen,%
    QDecorationStyled,%
    QSpacerItem,%
    QMdiSubWindow,%
    QFontComboBox,%
    QGraphicsBlurEffect,%
    QWSTslibMouseHandler,%
    QStyleOptionQ3ListViewItem,%
    QTimeEdit,%
    QGraphicsProxyWidget,%
    QTileRules,%
    QMatrix2x2,%
    QVFbKeyData,%
    QFontEngineFactoryInterface,%
    QIconEnginePluginV2,%
    QAccessibleBridgePlugin,%
    QStringListModel,%
    QStyleOptionSpinBox,%
    QStyleHintReturnMask,%
    QCheckBox,%
    QtEvents,%
    QWSTtyKeyboardHandler,%
    QGraphicsItem,%
    QPageSetupDialog,%
    QWSLinuxInputMouseHandler,%
    QMatrix4x3,%
    QWSSoundServer,%
    QGraphicsTextItem,%
    QGraphicsSceneResizeEvent,%
    QGraphicsScene,%
    QTextFrameLayoutData,%
    QDesktopWidget,%
    QTouchEvent,%
    QColorGroup,%
    QRadialGradient,%
    QWindowsMime,%
    QToolButton,%
    QPlastiqueStyle,%
    QButtonGroup,%
    QDecorationAction,%
    QPainterPathPrivate,%
    QDirModel,%
    QSortFilterProxyModel,%
    QWindowsVistaStyle,%
    QActionEvent,%
    QStyleOptionTabV2,%
    QStackedWidget,%
    QStatusBar,%
    QStyleOptionGroupBox,%
    QGraphicsSceneDragDropEvent,%
    QComboBox,%
    QStyleOptionSizeGrip,%
    QIconEngine,%
    QLinuxFbScreen,%
    QProxyStyle,%
    QKeyEvent,%
    QPictureFormatPlugin,%
    QFontDatabase,%
    QErrorMessage,%
    QInputMethodEvent,%
    QPanGesture,%
    QDataWidgetMapper,%
    QWidgetMapper,%
    QWSPcMouseHandler,%
    QBitmap,%
    QCursorShape,%
    QPrintPreviewDialog,%
    QAccessibleEditableTextInterface,%
    QAccessible,%
    QFontDialog,%
    QTextList,%
    QMouseEvent,%
    QStyleOptionTab,%
    QSpinBox,%
    QWSEmbedWidget,%
    QStyleHintReturn,%
    QGraphicsPixmapItem,%
    QAbstractButton,%
    QWindowsXPStyle,%
    QWSLinuxTPMouseHandler,%
    QImageTextKeyLang,%
    QListView,%
    QFontMetrics,%
    QSymbianEvent,%
    QUndoCommand,%
    QFontEnginePlugin,%
    QGraphicsScale,%
    QFileDialog,%
    QVBoxLayout,%
    QtDeclarative,%
    QDeclarativeExpression,%
    QDeclarativeListReference,%
    QDeclarativeAttachedPropertiesFunc,%
    QDeclarativeEngine,%
    QDeclarativeListProperty,%
    QDeclarativeError,%
    QDeclarativeComponent,%
    QDeclarativeScriptString,%
    QDeclarativeExtensionInterface,%
    QDeclarativeImageProvider,%
    QDeclarativePropertyValueInterceptor,%
    QDeclarativeProperties,%
    QDeclarativeContext,%
    QDeclarativeItem,%
    QDeclarativePropertyValueSource,%
    QDeclarativeView,%
    QDeclarativeTypeInfo,%
    QtDeclarative,%
    QDeclarativeExtensionPlugin,%
    QDeclarativeProperty,%
    QDeclarativeInfo,%
    QDeclarativePropertyMap,%
    QDeclarativeParserStatus,%
    QDeclarativeNetworkAccessManagerFactory,%
    QtTest,%
    QTestMouseEvent,%
    QTestLightXmlStreamer,%
    QTestFileLogger,%
    QTestData,%
    QTestXunitStreamer,%
    QSpontaneKeyEvent,%
    QTestEvent,%
    QTestEventList,%
    QEventSizeOfChecker,%
    QTestAccessibility,%
    QTestKeyClicksEvent,%
    QTestElementAttribute,%
    QTestAccessibilityEvent,%
    QTestCoreList,%
    QTestXmlStreamer,%
    QTestElement,%
    QTestBasicStreamer,%
    QSignalSpy,%
    QTestEventLoop,%
    QtTest,%
    QtTestGui,%
    QTestCoreElement,%
    QTest,%
    QTestKeyEvent,%
    QTestDelayEvent,%
    QtXml,%
    QDomNotation,%
    QDomImplementation,%
    QXmlDefaultHandler,%
    QXmlStreamReader,%
    QXmlInputSource,%
    QXmlStreamEntityDeclarations,%
    QXmlStreamWriter,%
    QXmlStreamAttributes,%
    QXmlReader,%
    QXmlNamespaceSupport,%
    QXmlDTDHandler,%
    QXmlLocator,%
    QDomNode,%
    QDomProcessingInstruction,%
    QXmlStreamNamespaceDeclaration,%
    QDomElement,%
    QXmlErrorHandler,%
    QDomDocumentFragment,%
    QXmlStreamEntityResolver,%
    QDomNamedNodeMap,%
    QDomAttr,%
    QXmlDeclHandler,%
    QDomEntity,%
    QXmlLexicalHandler,%
    QDomNodeList,%
    QDomCharacterData,%
    QXmlEntityResolver,%
    QXmlSimpleReader,%
    QXmlStreamAttribute,%
    QDomDocumentType,%
    QDomEntityReference,%
    QDomCDATASection,%
    QXmlStreamNotationDeclaration,%
    QXmlAttributes,%
    QXmlStreamStringRef,%
    QXmlStreamNotationDeclarations,%
    QtXml,%
    QDomComment,%
    QDomDocument,%
    QXmlContentHandler,%
    QXmlStreamNamespaceDeclarations,%
    QDomText,%
    QXmlParseException,%
    QXmlStreamEntityDeclaration,%
    QtWebKit,%
    QWebElement,%
    QWebSecurityOrigin,%
    QWebElementCollection,%
    QWebHistoryItem,%
    QGraphicsWebView,%
    QWebHistory,%
    QWebPage,%
    QWebFrame,%
    QWebSettings,%
    QWebHitTestResult,%
    QWebInspector,%
    QWebPluginFactory,%
    QWebDatabase,%
    QtWebKit,%
    QWebHistoryInterface,%
    QWebView,%
    Qt3Support,%
    Q3Button,%
    Q3SocketDevice,%
    Q3ProgressDialog,%
    Q3DataTable,%
    Q3IconDrag,%
    Q3HttpResponseHeader,%
    Q3GridLayout,%
    Q3ListBox,%
    Q3SqlFieldInfoList,%
    Q3DataBrowser,%
    Q3PaintDeviceMetrics,%
    Q3CanvasPolygon,%
    Q3PtrQueue,%
    Q3NetworkProtocolFactoryBase,%
    Q3SortedList,%
    Q3SingleCleanupHandler,%
    Q3NetworkOperation,%
    Q3PtrDictIterator,%
    Q3ObjectDictionary,%
    Q3CString,%
    Q3TextStream,%
    Q3DateTimeEdit,%
    Q3ButtonGroup,%
    Q3VBox,%
    Q3AsciiBucket,%
    Q3CanvasSprite,%
    Q3CanvasRectangle,%
    Q3HttpRequestHeader,%
    Q3UriDrag,%
    Q3CanvasPixmap,%
    Q3HBox,%
    Q3AsciiCache,%
    Q3GCacheIterator,%
    Q3GVector,%
    Q3ToolBar,%
    Q3NetworkProtocolFactory,%
    Q3GCache,%
    Q3CanvasText,%
    Q3IntDictIterator,%
    Q3VButtonGroup,%
    Q3TabDialog,%
    Q3UrlOperator,%
    Q3Action,%
    Q3TextEditOptimPrivate,%
    Q3SyntaxHighlighter,%
    Q3TableSelection,%
    Q3Http,%
    Q3GListStdIterator,%
    Q3GList,%
    Q3StrIList,%
    Q3PtrDict,%
    Q3PtrCollection,%
    Q3StyleSheetItem,%
    Q3GridView,%
    Q3ListViewItemIterator,%
    Q3LNode,%
    Q3Url,%
    Q3CanvasLine,%
    Q3ValueListConstIterator,%
    Q3FileDialog,%
    Q3EditorFactory,%
    Q3ComboBox,%
    Q3VBoxLayout,%
    Q3ProgressBar,%
    Q3Signal,%
    Q3HGroupBox,%
    Q3Picture,%
    Q3DragObject,%
    Q3PtrBucket,%
    Q3CanvasSpline,%
    Q3Accel,%
    Q3VGroupBox,%
    Q3DataView,%
    Q3CleanupHandler,%
    Q3TableItem,%
    Q3SqlRecordInfo,%
    Q3StrList,%
    Q3Ftp,%
    Q3IntCache,%
    Q3CheckListItem,%
    Q3StyleSheet,%
    Q3HBoxLayout,%
    Q3IntDict,%
    Q3FilePreview,%
    Q3DateTimeEditBase,%
    Q3Frame,%
    Q3AsciiDictIterator,%
    Q3GDict,%
    Q3ListViewItem,%
    Q3StoredDrag,%
    Q3StrListIterator,%
    Q3CanvasView,%
    Q3PtrStack,%
    Q3CacheIterator,%
    Q3SqlForm,%
    Q3RangeControl,%
    Q3Dns,%
    Q3Dict,%
    Q3ColorDrag,%
    Q3TextView,%
    Q3ValueListIterator,%
    Q3IntCacheIterator,%
    Q3Table,%
    Q3ComboTableItem,%
    Q3ScrollView,%
    Q3ListBoxPixmap,%
    Q3PopupMenu,%
    Q3SqlSelectCursor,%
    Q3DictIterator,%
    Q3ValueList,%
    Q3CanvasPolygonalItem,%
    Q3DateEdit,%
    Q3StrVec,%
    Q3GroupBox,%
    Q3DropSite,%
    Q3DnsSocket,%
    Q3PtrListStdIterator,%
    Q3WhatsThis,%
    Q3GDictIterator,%
    Q3DeepCopy,%
    Q3Cache,%
    Q3SimpleRichText,%
    Q3ListBoxText,%
    Q3IconViewItem,%
    Q3Grid,%
    Q3SpinWidget,%
    Q3BaseBucket,%
    Q3Painter,%
    Q3Header,%
    Q3NetworkProtocolDict,%
    Q3DockArea,%
    Q3Wizard,%
    Q3PtrList,%
    Q3MultiLineEdit,%
    Q3Shared,%
    Q3PointArray,%
    Q3ListBoxItem,%
    Q3ListView,%
    Q3ValueVector,%
    Q3WidgetStack,%
    Q3ServerSocket,%
    Q3MemArray,%
    Q3TimeEdit,%
    Q3AsciiCacheIterator,%
    Q3TSFUNC,%
    Q3HButtonGroup,%
    Q3CanvasItemList,%
    Q3AsciiDict,%
    Q3DockAreaLayout,%
    Q3IconDragItem,%
    Q3SqlEditorFactory,%
    Q3SqlPropertyMap,%
    Q3PtrVector,%
    Q3GListIterator,%
    Q3StrIVec,%
    Q3IntBucket,%
    Q3BoxLayout,%
    Q3CanvasItem,%
    Q3TextDrag,%
    Q3MainWindow,%
    Q3TextBrowser,%
    Q3GArray,%
    Q3PtrListIterator,%
    Q3CheckTableItem,%
    Q3HttpHeader,%
    Q3SqlFieldInfo,%
    Q3Socket,%
    Q3MimeSourceFactory,%
    Qt3Support,%
    Q3FileIconProvider,%
    Q3SqlCursor,%
    Q3StringBucket,%
    Q3CanvasPixmapArray,%
    Q3LocalFs,%
    Q3TextEdit,%
    Q3PolygonScanner,%
    Q3ActionGroup,%
    Q3NetworkProtocol,%
    Q3ValueStack,%
    Q3ImageDrag,%
    Q3Semaphore,%
    Q3CanvasEllipse,%
    Q3Canvas,%
    Q3DockWindow,%
    Q3IconView,%
    Q3Process,%
    QtXmlPatterns,%
    QXmlSerializer,%
    QXmlQuery,%
    QXmlNodeModelIndex,%
    QXmlSchema,%
    QXmlNamePool,%
    QXmlResultItems,%
    QSimpleXmlNodeModel,%
    QAbstractXmlNodeModel,%
    QSourceLocation,%
    QAbstractUriResolver,%
    QXmlSchemaValidator,%
    QXmlFormatter,%
    QXmlItem,%
    QXmlName,%
    QAbstractMessageHandler,%
    QtXmlPatterns,%
    QAbstractXmlReceiver,%
    QtHelp,%
    QHelpContentItem,%
    QHelpContentWidget,%
    QHelpEngine,%
    QHelpEngineCore,%
    QHelpSearchQuery,%
    QHelpIndexModel,%
    QHelpSearchResultWidget,%
    QHelpIndexWidget,%
    QHelpContentModel,%
    QHelpSearchEngine,%
    QHelpGlobal,%
    QHelpSearchQueryWidget,%
    QtHelp,%
    QtOpenVG,%
    QtOpenVG,%
    QtOpenGL,%
    QMacCompatGLenum,%
    QGLBuffer,%
    QGLPixelBuffer,%
    QGLColormap,%
    QMacCompatGLuint,%
    QGLFramebufferObjectFormat,%
    QGLScreenSurfaceFunctions,%
    QGLContext,%
    QGLFramebufferObject,%
    QGLScreen,%
    QGLWidget,%
    QMacGLCompatTypes,%
    QMacCompatGLint,%
    QGLShader,%
    QGLFormat,%
    QGLShaderProgram,%
    QtOpenGL%
  }
}


  \definecolor{devnet_magenta} {RGB}{255, 20,147} % {216,  0,117}
  \definecolor{devnet_green}   {RGB}{108,226,108} % { 58,216, 58}
  \definecolor{devnet_gray}    {RGB}{175,175,175} % {145,145,145}
  \definecolor{devnet_silver}  {RGB}{128,128,128} % { 98, 98, 98}
  \definecolor{devnet_blue}    {RGB}{  0,102,153} % {  0, 61, 92}
  \definecolor{devnet_fullblue}{RGB}{  0,  0,255} % {  0,  0,225}
  \definecolor{devnet_qtgreen} {RGB}{  0,130,  0} % {  0,100,  0}
                                                  % -30 lum points
  \lstset{ %
    language=[Qt]C++,
    basicstyle=\ttfamily\footnotesize,
    keywordstyle=\bfseries\color{devnet_blue},
    keywordstyle=[2]\bfseries\color{devnet_blue},
    keywordstyle=[3]\bfseries\color{devnet_magenta},
    stringstyle=\color{devnet_fullblue},
    commentstyle=\color{devnet_qtgreen},
    %numbers=left,
    %numberstyle=\tiny,
    %stepnumber=1,
    %numbersep=10pt,
    showspaces=false,
    showstringspaces=false,
    frame=l,
    framerule=1pt,
    framesep=4pt,
    %rulecolor=\color{black!30},
    rulecolor=\color{devnet_qtgreen},
    aboveskip=1mm,
    belowskip=1mm,
    showlines=false,
  }
}
{
  \lstdefinelanguage[Qt]{C++}[]{C++}%
  { keywordsprefix=Q
  }

  \lstset{ % 
    language=[Qt]C++,
    basicstyle=\ttfamily\footnotesize,
    keywordstyle=\bfseries,
    stringstyle=\color{blue},
    commentstyle=\color{qtlightgreen}, 
    %numbers=left,
    %numberstyle=\tiny,
    %stepnumber=1,
    %numbersep=10pt, 
    showspaces=false,
    showstringspaces=false,
    frame=l, 
    framerule=1pt,
    framesep=4pt,
    %rulecolor=\color{black!30},
    rulecolor=\color{qtlightgreen},
    aboveskip=1mm,
    belowskip=1mm,
    showlines=false,
  }
}

\lstnewenvironment{cpp}{\lstset{language=[Qt]C++}}{}
\lstnewenvironment{qmake}{\lstset{language=qmake}}{}
\lstnewenvironment{shell}{\lstset{language=bash}}{}
\lstnewenvironment{xml}{\lstset{language=html}}{}


\newcommand{\inputqml}[1]{%
\setlength{\lineskip}{0pt}
\footnotesize
\input{#1}
\normalsize
}

\newenvironment{qml}{%
\setlength{\lineskip}{0pt}
\footnotesize
}{%
}

\newcommand{\qthome}[2]{\href{http://qt.nokia.com/#1}{\beamerbutton{See #2}}}  
\newcommand{\doc}[2]{\href{http://qt.nokia.com/doc/latest/#1}{\beamerbutton{See #2 Documentation}}}  
\newcommand{\wiki}[2]{\href{http://en.wikipedia.org/wiki/#1}{\beamerbutton{See #2 Wikipedia Article}}}  
\newcommand{\doccreator}[2]{\href{http://doc.qt.nokia.com/qtcreator/#1}{\beamerbutton{See #2 Documentation}}}  
\newcommand{\demo}[1]{\href{run:#1}{\beamerbutton{Demo #1}}}  
\newcommand{\qtdemo}[1]{\href{run:#1}{\beamerbutton{Demo \$QTDIR/#1}}}  
\newcommand{\lab}[1]{\href{run:#1}{\beamerbutton{Lab #1}}}
\newcommand\externalLink[1]{{\sffamily\color{blue}{\href{#1}{#1}}}}
\newcommand\externalNamedLink[2]{{\sffamily\color{blue}{\href{#1}{#2}}}}


\date{}

\sloppy
\newenvironment{answers}[2]{%
\begin{frame}[fragile]
\begin{center}
#1\\(page #2)
\end{center}
\tiny
}{%
\end{frame}}

%----------------------------------------------------------------------

% setup for two column printing.
\setlength{\columnseprule}{1pt}
\setlength{\columnsep}{20pt}

%----------------------------------------------------------------------
\newcommand{\secTitle}{}
\newcommand{\setSectionTitle}[1]{\def\secTitle{#1}\section{#1}}
\newcommand{\sectionTitle}{\secTitle}

\newcommand{\contentHyperlink}[2]{\hyperlink{#2}{#1 (page~\pageref{#2})}}
\newcommand{\pageHyperlink}[1]{\hyperlink{#1}{page~\pageref{#1}}}
\newcommand{\itop}[2]{\item \contentHyperlink{#1}{#2}}
\newcommand{\itopStar}[2]{\item<*> \contentHyperlink{#1}{#2}}
\newcommand{\showHyperlink}[2]{\textcolor[rgb]{0,0,1}{\hyperlink{#1}{#2}}}

%----------------------------------------------------------------------
\newcommand{\qtThree}{$^{qt3}$}
\newcommand{\qt}[2]{$^{\mathrm{(Qt\,#1.#2)}}$}
\def\strikethrough#1{%
{%
  \setbox0=\hbox{#1}%
  \dimen0=\ht0 \dimen1=\dp0
  \setbox1=\vbox{
       \box0 \vskip-\dimen1\vskip-0.8ex \hrule
   }%
   \ht1=\dimen0 \dp1=\dimen1
   \box1
  }%
}

% These aren't really up to date currently
%\newcommand{\notQtopia}{$^{\tiny \hbox{\strikethrough{Qtopia}}}$}
%\newcommand{\Qtopia}{$^{\tiny \hbox{Qtopia}}$}
\newcommand{\notQtopia}{}
\newcommand{\Qtopia}{}
%----------------------------------------------------------------------
\newcommand{\pleaseNote}{\underline{\textbf{Note:}}\ }

%----------------------------------------------------------------------
%% All this is really magic to me, I found from pdfinfo slides.pdf that:
%% Page size:      362.835 x 272.126 pts
%% If I device 362.835 with the screen width, I get the magic
%% \paperToScreenFactor factor. Don't ask why or how I got to doing that!

%% Note: The \includegraphics command takes the resolution (dots-per-inch)
%% into account, so screenshots from higher resolution sources will be
%% scaled down even further.

\ifthenelse{\equal{\screenwidth}{1024}}{
  \newcommand{\paperToScreenFactor}{0.35433}
  \newcommand{\doublePaperToScreenFactor}{0.70866}
}{
  \ifthenelse{\equal{\screenwidth}{1280}}{
    \newcommand{\paperToScreenFactor}{0.28346}
    \newcommand{\doublePaperToScreenFactor}{0.56692}
  }{
  \ifthenelse{\equal{\screenwidth}{1400}}{
    \newcommand{\paperToScreenFactor}{0.25916}
    \newcommand{\doublePaperToScreenFactor}{0.51833}
    }{
    \error{Unknown resolution, acceptable screen width are 1024, 1280, and
      1400}
    }}
}

\newcommand{\flushedImage}[1]{\parpic[r]{\includegraphics[origin=t,scale=\paperToScreenFactor]{#1}}}
\newcommand{\flushedImageDoubleWidth}[1]{\parpic[r]{\includegraphics[origin=t,scale=\doublePaperToScreenFactor]{#1}}}
\newcommand{\image}[1]{\includegraphics[scale=\paperToScreenFactor]{#1}}
\newcommand{\imageDoubleWidth}[1]{\includegraphics[scale=\doublePaperToScreenFactor]{#1}}
\newcommand{\imageFullWidth}[1]{\includegraphics[width=0.75\pdfpagewidth,keepaspectratio]{#1}}
\newcommand{\centeredImage}[1]{\strut\hfill\includegraphics[scale=\paperToScreenFactor]{#1}\hfill\strut}
\newcommand{\centeredImageDoubleWidth}[1]{\strut\hfill\includegraphics[scale=\doublePaperToScreenFactor]{#1}\hfill\strut}
\newcommand{\centeredImageFullWidth}[1]{\strut\hfill\includegraphics[width=0.75\pdfpagewidth,keepaspectratio]{#1}\hfill\strut}
\newcommand{\centeredImageFullHeight}[1]{\strut\hfill\includegraphics[height=0.6\pdfpageheight,keepaspectratio]{#1}\hfill\strut}
\newcommand{\inlinedImage}[1]{\includegraphics[height=0.5cm]{#1}\strut}

% Open Office images do not need to be scale down, they are vector
% graphics, so we can use the same size for all resolutions.
\newcommand{\oooFlushedImage}[1]{\parpic[r]{\includegraphics[origin=t,scale=0.354331054688]{ooo/#1}}}
\newcommand{\oooImage}[1]{\includegraphics[scale=0.354331054688]{ooo/#1}}
\newcommand{\oooCenteredImage}[1]{\strut\hfill\includegraphics[scale=0.354331054688]{ooo/#1}\hfill\strut}

\definecolor{exampletitlecolor}{rgb}{0.9,0.9,1.0}
\newcommand{\exampletitle}[1]{\colorbox{exampletitlecolor}{\parbox{1.0\textwidth}{#1}}}


% Syntax coloring definitions
\definecolor{keyword}{rgb}{0.9,0.2,0.2}
\definecolor{class}{rgb}{0.0,0.0,0.9}
\definecolor{number}{rgb}{0.9,0.5,0.0}
\definecolor{type}{rgb}{0.6,0.6,0.0}
\definecolor{comment}{rgb}{0.2,0.5,0.2}
\definecolor{string}{rgb}{0.2,0.5,0.9}
\definecolor{operator}{rgb}{0.5,0.5,0.5}

% ----------- Notes (the boolean is set in setup.tex)
\ifthenelse{\boolean{with-notes}}{
  \includeversion{note}
}
{%else
  \excludeversion{note}
}
\excludeversion{EXCLUDE}

\newenvironment{instructors-notes}{%
\setbeamertemplate{headline} {\includegraphics[width=\textwidth]{images/comment-small}}
\setbeamertemplate{footline} {\includegraphics[width=\textwidth]{images/comment-small}}
    \begin{frame}[fragile,environment=instructors-notes]
}{%
\end{frame}}

% --------------- Watermark ---------------------------------------------
\ifthenelse{\boolean{eval-version}}{%
\let\evaltitle\frametitle
\renewcommand{\frametitle}[1]{\evaltitle{#1}%
  \begin{picture}(0,0)%
    \put(0,0){%
      \rotatebox{-25}{\textcolor[gray]{0.80}{\fontsize{15mm}{15mm}\selectfont Evaluation Copy}}}%
    \put(0,-20){%
      \rotatebox{-25}{\textcolor[gray]{0.70}{\fontsize{2mm}{2mm}\selectfont If you are being taught out of this training material or have been sold it, please contact qt-info@nokia.com}}}%
  \end{picture}%
}}{}


%% ------------ Title of the course ----------------------------------------
\newcommand{\frontpageTitleFormat}{\fontfamily{NokiaLarge}\fontseries{b}\fontsize{24pt}{28.8pt} \selectfont }

%------------ Setup of index command -----------------------
\newcommand{\classsep}{::}
\newcommand{\pointerStar}{*\xspace}
\newcommand{\programmingLanguage}{C++\xspace}
\newcommand{\pointerAmpersand}{\&}
\newcommand{\pointerDeref}{->}

%%% \newcommand{\capitalise}[1]{\capitalizefoo#1\end} \def\capitalizefoo#1#2\end{\uppercase{#1}\lowercase{#2}}
\newcommand{\possibleEmptyConcat}[2]{\ifthenelse{\equal{#1}{}}{#2}{#1 #2}}
\newcommand{\xConcept}[1]{\index{Concepts}{#1}}
\newcommand{\xCls}[1]{\index{Classes}{#1}}
\newcommand{\xNs}[1]{\index{Classes}{#1 (namespace)}}
\newcommand{\xNsCls}[2]{\index{Classes}{#1 (namespace)!#2}\index{Classes}{#2 (#1)}}
\newcommand{\xClsFn}[2]{\index{Classes}{#1!#2}\index{Classes}{#2 (#1)}}
\newcommand{\xNsFn}[2]{\index{Classes}{#1 (namespace)!#2}\index{Classes}{#2 (#1)}}
\newcommand{\xNsClsFn}[3]{\index{Classes}{#1 (namespace)!#2!#3}\index{Classes}{#2 (#1)!#3}\index{Classes}{#3 (#1\classsep#2)}}
%PENDING(blackie): have the extra section?
%\newcommand{\xClsSig}[2]{\index{Classes}{#1!#2 (signal)}\index{Classes}{#2 (#1, signal)}\index{Classes}{Signals!#1!#2}\index{Classes}{Signals!#2!#1}}
\newcommand{\xClsSig}[2]{\index{Classes}{#1!#2 (signal)}\index{Classes}{#2 (signal on #1)}}
\newcommand{\xNsClsSig}[3]{\index{Classes}{#1 (namespace)!#2!#3 (signal)}\index{Classes}{#2 (#1)!#3 (signal)}\index{Classes}{#3 (signal on #1\classsep#2)}}
%PENDING(blackie): have the extra section?
%\newcommand{\xClsSlt}[2]{\index{Classes}{#1!#2 (slot)}\index{Classes}{#2 (#1, slot)}\index{Classes}{Slots!#1!#2}\index{Classes}{Slots!#2!#1}}
\newcommand{\xClsSlt}[2]{\index{Classes}{#1!#2 (slot)}\index{Classes}{#2 (slot on #1)}}
\newcommand{\xNsClsSlt}[3]{\index{Classes}{#1 (namespace)!#2!#3 (slot)}\index{Classes}{#2 (#1)!#3 (slot)}\index{Classes}{#3 (slot on #1\classsep#2)}}
\newcommand{\xMacro}[1]{\index{Classes}{#1 (macro)}}
\newcommand{\xExample}[1]{\index{Examples}{#1}}
\newcommand{\xProject}[1]{\index{Projects}{#1}}

\newcommand{\iConcept}[2][]{#1{#2}\xConcept{#2}}
\newcommand{\iMacro}[2][\texttt]{#1{#2}\xMacro{#2}}
\newcommand{\iMacroPar}[3][\texttt]{#1{#2(#3)}\xMacro{#2}}
\newcommand{\iExample}[1]{\emph{#1}\xExample{#1}}
\newcommand{\iProject}[1]{\emph{#1}\xProject{#1}}

% \iCls*
\newcommand{\iCls}[2][\texttt]{#1{#2}\xCls{#2}}
\newcommand{\iClsFn}[3][\texttt]{#1{#2\classsep#3()}\xClsFn{#2}{#3}}
\newcommand{\iClsSig}[3][\texttt]{#1{#2\classsep#3()}\xClsSig{#2}{#3}}
\newcommand{\iClsSlt}[3][\texttt]{#1{#2\classsep#3()}\xClsSlt{#2}{#3}}
\newcommand{\iClsEnum}[3][\texttt]{#1{#2\classsep#3}\xClsFn{#2}{#3}}
\newcommand{\iClsFnPar}[4][\texttt]{#1{#2\classsep#3(#4)}\xClsFn{#2}{#3}}
\newcommand{\iClsSigPar}[4][\texttt]{#1{#2\classsep#3(#4)}\xClsSig{#2}{#3}}
\newcommand{\iClsSltPar}[4][\texttt]{#1{#2\classsep#3(#4)}\xClsSlt{#2}{#3}}

% \iClsT*
\newcommand{\iClsT}[2][T]{\texttt{#2<#1>}\xCls{#2}}
\newcommand{\iClsTFn}[3][T]{\texttt{#2<#1>\classsep#3()}\xClsFn{#2}{#3}}
\newcommand{\iClsTSig}[3][T]{\texttt{#2<#1>\classsep#3()}\xClsSig{#2}{#3}}
\newcommand{\iClsTSlt}[3][T]{\texttt{#2<#1>\classsep#3()}\xClsSlt{#2}{#3}}
\newcommand{\iClsTEnum}[3][T]{\texttt{#2<#1>\classsep#3}\xClsFn{#2}{#3}}
\newcommand{\iClsTFnPar}[4][T]{\texttt{#2<#1>\classsep#3(#4)}\xClsFn{#2}{#3}}
\newcommand{\iClsTSigPar}[4][T]{\texttt{#2<#1>\classsep#3(#4)}\xClsSig{#2}{#3}}
\newcommand{\iClsTSltPar}[4][T]{\texttt{#2<#1>\classsep#3(#4)}\xClsSlt{#2}{#3}}

% \iNs*
\newcommand{\iNs}[2][\texttt]{#1{#2}\xNs{#2}}
\newcommand{\iNsFn}[3][\texttt]{#1{#2\classsep#3()}\xNsFn{#2}{#3}}
\newcommand{\iNsEnum}[3][\texttt]{#1{#2\classsep#3}\xNsFn{#2}{#3}}
\newcommand{\iNsFnPar}[4][\texttt]{#1{#2\classsep#3(#4)}\xNsFn{#2}{#3}}

% \iNsCls*
\newcommand{\iNsCls}[3][\texttt]{#1{#2\classsep#3}\xNsCls{#2}{#3}}
\newcommand{\iNsClsFn}[4][\texttt]{#1{#2\classsep#3\classsep#4()}\xNsClsFn{#2}{#3}{#4}}
\newcommand{\iNsClsSig}[4][\texttt]{#1{#2\classsep#3\classsep#4()}\xNsClsSig{#2}{#3}{#4}}
\newcommand{\iNsClsSlt}[4][\texttt]{#1{#2\classsep#3\classsep#4()}\xNsClsSlt{#2}{#3}{#4}}
\newcommand{\iNsClsEnum}[4][\texttt]{#1{#2\classsep#3\classsep#4}\xNsClsFn{#2}{#3}{#4}}
\newcommand{\iNsClsFnPar}[5][\texttt]{#1{#2\classsep#3\classsep#4(#5)}\xNsClsFn{#2}{#3}{#4}}
\newcommand{\iNsClsSigPar}[5][\texttt]{#1{#2\classsep#3\classsep#4(#5)}\xClsSig{#2}{#3}{#4}}
\newcommand{\iNsClsSltPar}[5][\texttt]{#1{#2\classsep#3\classsep#4(#5)}\xClsSlt{#2}{#3}{#4}}

% like \iCls*, but hide the class:
\newcommand{\hClsFn}[3][\texttt]{#1{#3()}\xClsFn{#2}{#3}}
\newcommand{\hClsSig}[3][\texttt]{#1{#3()}\xClsSig{#2}{#3}}
\newcommand{\hClsSlt}[3][\texttt]{#1{#3()}\xClsSlt{#2}{#3}}
\newcommand{\hClsEnum}[3][\texttt]{#1{#3}\xClsFn{#2}{#3}}
\newcommand{\hClsFnPar}[4][\texttt]{#1{#3(#4)}\xClsFn{#2}{#3}}
\newcommand{\hClsSigPar}[4][\texttt]{#1{#3(#4)}\xClsSig{#2}{#3}}
\newcommand{\hClsSltPar}[4][\texttt]{#1{#3(#4)}\xClsSlt{#2}{#3}}

% like \iNs*, but hide the namespace:
\newcommand{\hNsFn}[3][\texttt]{#1{#3()}\xNsFn{#2}{#3}}
\newcommand{\hNsEnum}[3][\texttt]{#1{#3}\xNsFn{#2}{#3}}
\newcommand{\hNsFnPar}[4][\texttt]{#1{#3(#4)}\xNsFn{#2}{#3}}

% like \iNsCls*, but hide namespace (but not class):
\newcommand{\hNsCls}[3][\texttt]{#1{#3}\xNsCls{#2}{#3}}
\newcommand{\hNsClsFn}[4][\texttt]{#1{#3\classsep#4()}\xNsClsFn{#2}{#3}{#4}}
\newcommand{\hNsClsSig}[4][\texttt]{#1{#3\classsep#4()}\xNsClsSig{#2}{#3}{#4}}
\newcommand{\hNsClsSlt}[4][\texttt]{#1{#3\classsep#4()}\xNsClsSlt{#2}{#3}{#4}}
\newcommand{\hNsClsEnum}[4][\texttt]{#1{#3\classsep#4}\xNsClsFn{#2}{#3}{#4}}
\newcommand{\hNsClsFnPar}[5][\texttt]{#1{#3\classsep#4(#5)}\xNsClsFn{#2}{#3}{#4}}
\newcommand{\hNsClsSigPar}[5][\texttt]{#1{#3\classsep#4(#5)}\xNsClsSig{#2}{#3}{#4}}
\newcommand{\hNsClsSltPar}[5][\texttt]{#1{#3\classsep#4(#5)}\xNsClsSlt{#2}{#3}{#4}}

% like \iNsCls*, but hide namespace /and/ class:
\newcommand{\hhNsClsFn}[4][\texttt]{#1{#4()}\xNsClsFn{#2}{#3}{#4}}
\newcommand{\hhNsClsSig}[4][\texttt]{#1{#4()}\xNsClsSig{#2}{#3}{#4}}
\newcommand{\hhNsClsSlt}[4][\texttt]{#1{#4()}\xNsClsSlt{#2}{#3}{#4}}
\newcommand{\hhNsClsEnum}[4][\texttt]{#1{#4}\xNsClsFn{#2}{#3}{#4}}
\newcommand{\hhNsClsFnPar}[5][\texttt]{#1{#4(#5)}\xNsClsFn{#2}{#3}{#4}}
\newcommand{\hhNsClsSigPar}[5][\texttt]{#1{#4(#5)}\xNsClsSig{#2}{#3}{#4}}
\newcommand{\hhNsClsSltPar}[5][\texttt]{#1{#4(#5)}\xNsClsSlt{#2}{#3}{#4}}

\newcommand{\iClsVar}{\iClsEnum}
\newcommand{\hClsVar}{\hClsEnum}
\newcommand{\iClsTypedef}{\iClsEnum}
\newcommand{\hClsTypedef}{\hClsEnum}

\newcommand{\iNsVar}{\iNsEnum}
\newcommand{\hNsVar}{\hNsEnum}
\newcommand{\iNsTypedef}{\iNsEnum}
\newcommand{\hNsTypedef}{\hNsEnum}

\newcommand{\iNsClsVar}{\iNsClsEnum}
\newcommand{\hNsClsVar}{\hNsClsEnum}
\newcommand{\hhNsClsVar}{\hhNsClsEnum}
\newcommand{\iNsClsTypedef}{\iNsClsEnum}
\newcommand{\hNsClsTypedef}{\hNsClsEnum}
\newcommand{\hhNsClsTypedef}{\hhNsClsEnum}

% expands to "example $QTDIR/examples/#1"
\newcommand{\qtexample}[1]{example \emph{\$QTDIR/examples/\iExample{#1}}}
% expands to "handout/#1"
\newcommand{\projecthandout}[1]{\emph{handout/#1}}
% expands to "solutions/#1"
\newcommand{\projectsolution}[1]{\emph{solutions/#1}}

%-------------------------- Review commands--------------------------------------------
\newcommand{\reviewTT}{%
  \begin{picture}(0,0)%
    \put(313,40){%
      \rotatebox{-90}{\textcolor{red}{\fontsize{10mm}{10mm}\selectfont
          Qt REVIEW}}%
      }%
  \end{picture}%
  \par
 \vspace{-6mm}
}

\newboolean{ignore-slide}
\excludeversion{EXCLUDE}
%--------------------------- Slide environment -------------------------------------------
\newenvironment{slide}[2][1]{%
  \renewcommand{\slideId}{#2}
  \ifthenelse{\equal{#1}{1}}{
    \begin{frame}[environment=slide]
      }{
    \begin{frame}[#1,environment=slide]
  }
}%
{\end{frame}}

\AtBeginSection[] % Do nothing for \section*
{
  \begin{frame}<beamer>  
    \frametitle{Module: \secname}
     \tableofcontents[sectionstyle=hide,subsectionstyle=show/show/hide]
 \end{frame}
}

\AtBeginSubsection[] % Do nothing for \subsection*
{
  \begin{frame}<beamer>
    \frametitle{Module: \secname}
    \medskip
    \tableofcontents[sectionstyle=hide,subsectionstyle=show/shaded/hide]  
  \end{frame}
}

\AtBeginPart
{
  \begin{frame}<beamer>
    \frametitle{\insertshortpart}
    \tableofcontents[sectionstyle=show,subsectionstyle=hide]      
  \end{frame}
}


\newcommand{\todo}[1]{\begin{minipage}{3cm}
        \tiny{#1} \end{minipage}}


\newenvironment{questionize}{\begin{itemize}[<+>]}{\end{itemize}}

\ifthenelse{\boolean{cc-license}}
{
  \usebackgroundtemplate{\includegraphics[width=\paperwidth]{images/background_cc}}
} % else
{
  \usebackgroundtemplate{\includegraphics[width=\paperwidth]{images/background}}
}

% Highlighting in {alltt}, slides style:
\newcommand{\highlight}[1]{\textcolor{red}{#1}}

\newcommand{\correct}{\textcolor{qtdarkgreen}{$\surd$}}
\newcommand{\incorrect}{\textcolor{red}{$X$}}
\newcommand{\incorrecttext}[1]{\textcolor{red}{\strikethrough{#1}}}

\newcommand{\ok}{\textcolor{qtdarkgreen}{$OK$}}
\newcommand{\new}{\textcolor{qtdarkgreen}{$NEW$}}
\newcommand{\notok}{\textcolor{red}{$NOK$}}
\newcommand{\work}{\textcolor{yellow}{$WORK$}}
\newcommand{\unknown}{\textcolor{orange}{$???$}}
\newcommand{\delete}{\textcolor{red}{$DEL$}}

\newcommand{\consolecode}[1]{%
\vspace*{0.25em}
\fcolorbox{gray}{darkgray}{\begin{minipage}{\columnwidth - \leftmargin}
\textcolor{white}{\texttt{#1}}
\end{minipage}}
\vspace*{0.25em}}

\newcommand{\qc}[2]{\textcolor{#1}{#2}}
\newcommand{\qic}[2]{{\small\texttt{\textcolor{#1}{#2}}}}
\newcommand{\qtt}[1]{\texttt{#1}}

% cc-license fixed configuration. Automatically applied if in config.tex cc-license is set to true

\ifthenelse{\boolean{cc-license}}
  {
  }
  {
  % \customer{}
  % \partner{\includegraphics[height=2cm]{images/qt-logo} \\ http://qt.nokia.com}
  }

\usepackage{fontspec}
\usepackage{xunicode}

% set the default roman font
\setmainfont[Scale=0.9]{Arial}

% set the default sans font
\setsansfont[Scale=0.9]{Arial}

% set the default mono font
\setmonofont[Scale=0.8]{Liberation Mono}

% define new font families for main body, footnote, title
\newfontfamily\slidebodyfont{Arial}

\newfontfamily\slidefootnotefont{Arial}

\newfontfamily\slidetitlefont{Arial Bold}



\coursetitle{Qt Quick for Qt Developers}

\begin{document}
%%%%%%%%%%%%%%%%%%%%%%%%%%%%%%%%%%%%%%%%%%%%%%%%%%%%%%%%%%%%%%%%%%%%%%%%%%
%
% Copyright (c) 2008-2011, Nokia Corporation and/or its subsidiary(-ies).
% All rights reserved.
%
% This work, unless otherwise expressly stated, is licensed under a
% Creative Commons Attribution-ShareAlike 2.5.
%
% The full license document is available from
% http://creativecommons.org/licenses/by-sa/2.5/legalcode .
%
%%%%%%%%%%%%%%%%%%%%%%%%%%%%%%%%%%%%%%%%%%%%%%%%%%%%%%%%%%%%%%%%%%%%%%%%%%

% background for title page
\ifthenelse{\boolean{cc-license}}
{
  \setbeamertemplate{background canvas}{\includegraphics[width=\paperwidth,height=\paperheight]{images/titlepage-background_cc}}
  \begin{slide}[plain]{0449}
    \titlepage
  \end{slide}

  % simple slide background
  \setbeamertemplate{background canvas}{\includegraphics[width=\paperwidth,height=\paperheight]{images/background_cc}}
} % else
{
  \setbeamertemplate{background canvas}{\includegraphics[width=\paperwidth,height=\paperheight]{images/titlepage-background}}
  \begin{slide}[plain]{0449}
    \titlepage
  \end{slide}

  % simple slide background
  \setbeamertemplate{background canvas}{\includegraphics[width=\paperwidth,height=\paperheight]{images/background}}
}


\section{Introduction to Qt Quick}

%----------------------------------------------------------------------
\begin{slide}{2024}\frametitle{Objectives}

\begin{itemize}
\item Overview of the Qt library
  \begin{itemize}
  \item Qt framework presentation
  \item Qt Quick inside the Qt framework
  \end{itemize}\medskip
\item Understanding of QML syntax and concepts
  \begin{itemize}
  \item elements and identities
  \item properties and property binding
  \end{itemize}\medskip
\item Basic user interface composition skills
  \begin{itemize}
  \item familiarity with common elements
  \item understanding of anchors and their uses
  \item ability to reproduce a design
  \end{itemize}
\end{itemize}

\end{slide}

%%%%%%%%%%%%%%%%%%%%%%%%%%%%%%%%%%%%%%%%%%%%%%%%%%%%%%%%%%%%%%%%%%%%%%%%%%
%
% Copyright (c) 2008-2011, Nokia Corporation and/or its subsidiary(-ies).
% All rights reserved.
%
% This work, unless otherwise expressly stated, is licensed under a
% Creative Commons Attribution-ShareAlike 2.5.
%
% The full license document is available from
% http://creativecommons.org/licenses/by-sa/2.5/legalcode .
%
%%%%%%%%%%%%%%%%%%%%%%%%%%%%%%%%%%%%%%%%%%%%%%%%%%%%%%%%%%%%%%%%%%%%%%%%%%

%----------------------------------------------------------------------
\subsection{Meet Qt}

%----------------------------------------------------------------------
\begin{slide}{2003}\frametitle{The Qt SDK}
\vspace*{0.5em}
\centeredImageFullWidth{qml-intro/images/qt_architecture.png}

\end{slide}

%----------------------------------------------------------------------
\begin{slide}{2002}\frametitle{Qt modules}
\vspace*{1.5em}
The Qt framework is split in modules:
\begin{itemize}
\item Each module has is own repository
\item QtBase is the only mandatory module for all Qt application
\item QtTools is a special module containing Qt applications
\item examples: QtDeclarative, QtWebKit, QtMultimedia...
\end{itemize}\medskip
Modules are themselves split into one or more libraries
\begin{itemize}
\item libraries are linked to your applications
\item they group a set of common features (xml, dbus, network...)
\item QtCore (part of QtBase) is mandatory for all Qt applications
\item examples: QtNetwork, QtGui, QtDeclrative
\end{itemize}

\end{slide}

%----------------------------------------------------------------------

\begin{slide}{2001}\frametitle{Creating user interface}
\vspace*{1.5em}
Using widgets (C++)
\begin{itemize}
\item fixed interface
\item mostly targeting desktop application
\end{itemize}
\centeredImageFullWidth{qml-intro/images/ui-widget.png}

\end{slide}

%----------------------------------------------------------------------

\begin{slide}{2001}\frametitle{Creating user interface}
\vspace*{1.5em}
Using graphics view (C++)
\begin{itemize}
\item 2D canvas (with z axis)
\item handle a high number of 2D items
\item transformations and animations of items
\end{itemize}
\centeredImageFullWidth{qml-intro/images/ui-graphicsview.png}

\end{slide}

%----------------------------------------------------------------------

\begin{slide}{2001}\frametitle{Creating user interface}
\vspace*{1.5em}
Using openGL (C++/OpenGL)
\begin{itemize}
\item full strength of OpenGL
\item user interaction (mouse and keybboard) using Qt API
\item all classes from Qt are usable
\end{itemize}
%\centeredImageFullWidth{qml-intro/images/ui-opengl.png}

\end{slide}

%----------------------------------------------------------------------

\begin{slide}{2001}\frametitle{Creating user interface}
\vspace*{1.5em}
Using QtQuick (QML/Javascript/C++/OpenGL)
\begin{itemize}
\item declarative langage using the QML langage
\item provide smooth animation
\item easily linked to C++
\end{itemize}
\centeredImageFullWidth{qml-intro/images/ui-quick.png}

\end{slide}

%----------------------------------------------------------------------

\begin{slide}{2000}\frametitle{Qt Quick requirement}
\vspace*{1.5em}
\begin{itemize}
\item Platform must support OpenGL ES2
\item Needs at least QtBase, QtJsBackend and QtDeclarative module
\item Other module can be used to add new features:
\begin{itemize}
\item QtGraphicalEffects: add effects like blur, drop shadow...
\item Qt3D: 3D programming in QML
\item QtMultimedia: audio and video items
\item ...
\end{itemize}
\end{itemize}

\end{slide}

%----------------------------------------------------------------------

%%%%%%%%%%%%%%%%%%%%%%%%%%%%%%%%%%%%%%%%%%%%%%%%%%%%%%%%%%%%%%%%%%%%%%%%%%
%
% Copyright (c) 2008-2011, Nokia Corporation and/or its subsidiary(-ies).
% All rights reserved.
%
% This work, unless otherwise expressly stated, is licensed under a
% Creative Commons Attribution-ShareAlike 2.5.
%
% The full license document is available from
% http://creativecommons.org/licenses/by-sa/2.5/legalcode .
%
%%%%%%%%%%%%%%%%%%%%%%%%%%%%%%%%%%%%%%%%%%%%%%%%%%%%%%%%%%%%%%%%%%%%%%%%%%

%----------------------------------------------------------------------
\subsection{Meet Qt Quick}

%----------------------------------------------------------------------
\begin{slide}{2003}\frametitle{What is Qt Quick?}
\vspace*{1.5em}

A set of technologies including:
\begin{itemize}
\item Declarative markup language: QML
\item Language runtime integrated with Qt
\item Qt Creator IDE support for the QML language
\item Graphical design tool
\item C++ API for integration with Qt applications
\end{itemize}
\end{slide}

%----------------------------------------------------------------------
\begin{slide}{2002}\frametitle{Philosophy of Qt Quick}
\vspace*{1.5em}

\begin{itemize}
\item Intuitive User Interfaces
\item Design-Oriented
\item Rapid Prototyping and Production
\item Easy Deployment
\end{itemize}
\end{slide}

%----------------------------------------------------------------------

\begin{slide}{2001}\frametitle{Design Criteria: Intuitive User Interfaces}

% Pictures of new-style user interfaces


\begin{itemize}
\item More natural ways to interact with applications
\item Easier to predict what elements will do
\item Smooth motion like real-world objects
\end{itemize}

\end{slide}

%----------------------------------------------------------------------

\begin{slide}{2000}\frametitle{Design Criteria: Design-Oriented}

% Pictures of Qt Creator QML integration

Traditional Qt way:
\begin{itemize}
\item Design tools for developers
\item Rich, platform-native widgets
\end{itemize}

% Qt way to QML way (rich widgets -> poor elements)

The Qt Quick way:
\begin{itemize}
\item Development tools and processes that are accessible to designers
\item Lightweight, customizable elements
\end{itemize}

\end{slide}

%----------------------------------------------------------------------

\begin{slide}{1999}\frametitle{Design Criteria: Rapid Prototyping and Production}

\begin{itemize}
    \item No compilation step
    \item JavaScript used as scripting language
    \begin{itemize}
       \item See \emph{JavaScript: The Definitive Guide} 
       \item See \externalLink{https://developer.mozilla.org/en/JavaScript}
    \end{itemize}
    \item Helpful to understand:
    \begin{itemize}
       \item HTML, CSS 
       \item ... but not required  
    \end{itemize} 
    \item Requires little or no programming experience
\end{itemize}

% JavaScript is a prototype-based programming language
% - slightly different paradigm to object-oriented programming
% http://www.ecma-international.org/publications/standards/Ecma-262.htm

\end{slide}

%----------------------------------------------------------------------

\begin{slide}{1998}\frametitle{Design Criteria: Easy Deployment}

\begin{itemize}
\item Self-contained packages
\item Installation is optional
\item Network transparent
  \begin{itemize}
  \item allows deployment over a network
  \item rich client front end to online services
  \end{itemize}
\end{itemize}

\end{slide}

%----------------------------------------------------------------------

%%%%%%%%%%%%%%%%%%%%%%%%%%%%%%%%%%%%%%%%%%%%%%%%%%%%%%%%%%%%%%%%%%%%%%%%%%
%
% Copyright (c) 2008-2011, Nokia Corporation and/or its subsidiary(-ies).
% All rights reserved.
%
% This work, unless otherwise expressly stated, is licensed under a
% Creative Commons Attribution-ShareAlike 2.5.
%
% The full license document is available from
% http://creativecommons.org/licenses/by-sa/2.5/legalcode .
%
%%%%%%%%%%%%%%%%%%%%%%%%%%%%%%%%%%%%%%%%%%%%%%%%%%%%%%%%%%%%%%%%%%%%%%%%%%

%----------------------------------------------------------------------

\subsection{Concepts}
\begin{slide}{2023}\frametitle{What is QML?}
\vspace*{0.5em}

% Basic concepts - what is Qt Quick, declarative UI?
Declarative language for User Interface elements:

\vspace*{0.5em}
\begin{itemize}
    \item Describes the user interface
    \begin{itemize}
    \item What elements look like
    \item How elements behave
    \end{itemize}
\item UI specified as tree of elements with properties     
\end{itemize}

\end{slide}

%----------------------------------------------------------------------

\begin{slide}{2022}\frametitle{A Tree of Elements}

% Insert tree <-> UI comparison image here.

% Other than the fundamental difference inheritance/parent-child trees that
% needs to be communicated, is there some fundamental difference between
% nested elements and elements listed using the children property.
% (Probably an implementation detail we don't want to address.)
                               
\centeredImageDoubleWidth{qml-intro/images/tree-of-elements.pdf}

\vspace*{1em}
Let's start with an example...

\end{slide}

%----------------------------------------------------------------------

\begin{slide}[fragile]{2021}\frametitle{Viewing an Example}

% declarative-uis/concepts/rectangle.qml
\inputqml{qml-intro/colorized/rectangle}

% Launching an example - the QML viewer, some simple markup
\vspace*{1em}
\begin{itemize}
\item Locate the example: \texttt{rectangle.qml}
\item Launch the QML viewer:\\
\consolecode{qmlscene~rectangle.qml}

% qml(viewer) shows the top level element at the correct size initially, but
% it will be resized with the window.
%
% Default values for x and y.

%%% Explain QML viewer                  

\demo{qml-intro/ex-concepts/rectangle.qml}

\end{itemize}

\end{slide}

%----------------------------------------------------------------------

\begin{slide}{2020}\frametitle{Example Concepts}

\begin{qml}
\qtt{\qc{keyword}{import}~\qc{class}{QtQuick}~\qc{number}{2.0}}\\
\vspace*{0.5em}
\qtt{\qc{lightgray}{// Define a light blue square}}\\
\vspace*{0.5em}
\qtt{\qc{lightgray}{Rectangle~\{}}\\
\qtt{~~~~\qc{lightgray}{width:~400;~height:~400}}\\
\qtt{~~~~\qc{lightgray}{color:~"lightblue"}}\\
\qtt{\qc{lightgray}{\}}}
\end{qml}

\vspace*{1em}
\begin{itemize}
\item To use Qt's features, \qic{keyword}{import} the \qic{class}{Qt} module
\item Specify a version to get the features you want
  \begin{itemize}
  \item only imports the features from that version
  \item will not use features from later versions, even if available
  \item will not fall back to features from an earlier version
  \end{itemize}
\item Guarantees that the behavior of the code will not change
  \begin{itemize}
  \item modules include support for multiple versions
  \item an upgraded module retains support for older versions of itself
  \end{itemize}
\end{itemize}

%%% Imports of other versions

\end{slide}

%----------------------------------------------------------------------

\begin{slide}{2019}\frametitle{Example Concepts}

\begin{qml}
\qtt{\qc{keyword}{import}~\qc{class}{QtQuick}~\qc{number}{2.0}}\\
\vspace*{0.5em}
\qtt{\qc{comment}{// Define a light blue square}}\\
\vspace*{0.5em}
\qtt{\qc{lightgray}{Rectangle~\{}}\\
\qtt{~~~~\qc{lightgray}{width:~400;~height:~400}}\\
\qtt{~~~~\qc{lightgray}{color:~"lightblue"}}\\
\qtt{\qc{lightgray}{\}}}
\end{qml}

\vspace*{1em}
\begin{itemize}
\item Use \qic{comment}{//} to add single line comments
\item Put multi-line comments inside \qic{comment}{/*} and \qic{comment}{*/}
\end{itemize}

\end{slide}

%----------------------------------------------------------------------

\begin{slide}{2018}\frametitle{Example Concepts}

\begin{qml}
\qtt{\qc{lightgray}{import~QtQuick~2.0}}\\
\vspace*{0.5em}
\qtt{\qc{lightgray}{// Define a light blue square}}\\
\vspace*{0.5em}
\qtt{\qc{class}{Rectangle}~\{}\\
\qtt{~~~~\qc{lightgray}{width:~400;~height:~400}}\\
\qtt{~~~~\qc{lightgray}{color:~"lightblue"}}\\
\qtt{\}}
\end{qml}

\vspace*{1em}
\begin{itemize}
\item Declare the elements you want to use
\item Each element has a body between \texttt{\{} and \texttt{\}}
\item A set of default elements are included in the \texttt{\qc{class}{Qt}} module
\end{itemize}

\end{slide}

%----------------------------------------------------------------------

\begin{slide}{2017}\frametitle{Example Concepts}

\begin{qml}
\qtt{\qc{lightgray}{import~QtQuick~2.0}}\\
\vspace*{0.5em}
\qtt{\qc{lightgray}{// Define a light blue square}}\\
\vspace*{0.5em}
\qtt{\qc{lightgray}{Rectangle~\{}}\\
\qtt{~~~~\qc{type}{width}:~\qc{number}{400};~\qc{type}{height}:~\qc{number}{400}}\\
\qtt{~~~~\qc{type}{color}:~\qc{string}{"lightblue"}}\\
\qtt{\qc{lightgray}{\}}}
\end{qml}

\vspace*{1em}
\begin{itemize}
\item Elements contain properties
\item Each property is defined using its name and a value
\item \texttt{\qc{type}{name}~:~value}
\end{itemize}

\end{slide}

%----------------------------------------------------------------------

\begin{slide}{2016}\frametitle{Example Summary}

\flushedImage{qml-intro/images/rectangle.pdf}
\begin{qml}
\qtt{\qc{keyword}{import}~\qc{class}{QtQuick}~\qc{number}{2.0}}\\
\vspace*{0.5em}
\qtt{\qc{comment}{// Define a light blue square}}\\
\vspace*{0.5em}
\qtt{\qc{class}{Rectangle}~\{}\\
\qtt{~~~~\qc{type}{width}:~\qc{number}{400};~\qc{type}{height}:~\qc{number}{400}}\\
\qtt{~~~~\qc{type}{color}:~\qc{string}{"lightblue"}}\\
\qtt{\}}
\end{qml}

\vspace*{1em}
\begin{itemize}
\item A \texttt{\textcolor{class}{Rectangle}} element
  \begin{itemize}
  \item with a body: \texttt{\{ ... \}}
  \item containing \qic{type}{width}, \qic{type}{height} and \qic{type}{color}
  properties
  \item separated by line breaks or semicolons
  \end{itemize}
\item Running the example in the \texttt{qml} viewer
  \begin{itemize}
  \item the viewer window will be 400 by 400
  \end{itemize}
\end{itemize}

\end{slide}

%----------------------------------------------------------------------

\begin{slide}{2015}\frametitle{Elements}

\begin{itemize}
\item Elements are structures in the markup language
  \begin{itemize}
  \item represent visible and non-visible parts
  \end{itemize}
\item \qic{class}{Item} is the base type of visible elements
  \begin{itemize}
  \item not visible itself
  \item has a position, dimensions
  \item usually used to group visual elements
  \item often used as the top-level element
  \item \qic{class}{Rectangle}, \qic{class}{Text}, \qic{class}{TextInput}, ...
  \end{itemize}
\item Non-visible elements:
  \begin{itemize}
  \item states, transitions, ...
  \item models, paths, ...
  \item gradients, timers, etc.
  \end{itemize}
\item Elements contain properties
  \begin{itemize}
  \item can also be extended with custom properties
  \end{itemize}
\end{itemize}

% A full list of elements and items is available in the documentation.
% Alternatively, you can look in the Qt Creator sources for the XML
% file that defines the basic types.

%%% Items based on QGraphicsItems
%%% qml viewer uses QGraphicsView
%%% Element hierarchy

\doc{qdeclarativeelements.html}{QML Elements}
\end{slide}

%----------------------------------------------------------------------

\begin{slide}{2014}\frametitle{Properties}

Elements are described by properties:

\begin{itemize}
\item Simple name-value definitions
  \begin{itemize}
  \item \qic{type}{width}, \qic{type}{height}, \qic{type}{color}, ...
  \item with default values
  \item each has a well-defined type
  \item separated by semicolons or line breaks
  \end{itemize}
  \vspace*{0.5em}
\item Used for
  \begin{itemize}
  \item identifying elements (\qic{type}{id} property)
  \item customizing their appearance
  \item changing their behavior
  \end{itemize}
  \vspace*{0.5em}
\end{itemize}

\end{slide}

%----------------------------------------------------------------------

\begin{slide}{2013}\frametitle{Property Examples}

\begin{itemize}
\item \textbf{Standard properties} can be given values:

\vspace*{0.25em}
\inputqml{qml-intro/colorized/standard-properties}

\vspace*{0.5em}
\item \textbf{Grouped properties} keep related properties together:

\vspace*{0.25em}
\inputqml{qml-intro/colorized/grouped-properties}

\vspace*{0.5em}
\item \textbf{Identity property} gives the element a name:

\vspace*{0.25em}
\inputqml{qml-intro/colorized/id-property}
\end{itemize}

\end{slide}

%----------------------------------------------------------------------

\begin{slide}{2012}\frametitle{Property Examples}

\begin{itemize}
\item \textbf{Attached properties} are applied to elements:

\vspace*{0.25em}
\inputqml{qml-intro/colorized/attached-properties}

\begin{itemize}
\item \qic{type}{KeyNagivation.tab} is not a standard property of \qic{class}{TextInput}
\item is a standard property that is attached to elements
\end{itemize}

\vspace*{0.5em}
\item \textbf{Custom properties} can be added to any element:

\vspace*{0.25em}
\inputqml{qml-intro/colorized/custom-property}
\end{itemize}

\end{slide}

%%% Add examples of properties with some standard elements
%%% Mention attached properties, but explain them when showing key input
%%% Highlight custom properties, but explain them 

%----------------------------------------------------------------------

\begin{slide}{2011}\frametitle{Binding Properties}

\flushedImage{qml-intro/images/expressions.png}
% declarative-uis/concepts/expressions.qml
\inputqml{qml-intro/colorized/expressions}

\demo{qml-intro/ex-concepts/expressions.qml}

\begin{itemize}
\item Properties can contain expressions
  \begin{itemize}
  \item see above: \qic{type}{width} is twice the \qic{type}{height}
  \end{itemize}
\item Not just initial assignments
\item Expressions are evaluated when needed
\end{itemize}

\doc{propertybinding.html}{Property Binding}
\end{slide}

%----------------------------------------------------------------------

\begin{slide}{2010}\frametitle{Identifying Elements}

% id is an especially important property because it provides a name
% for an element
% Coding convention - always first property?

The \qic{type}{id} property defines an identity for an element

\begin{itemize}
\item Lets other elements refer to it
  \begin{itemize}
  \item for relative alignment and positioning
  \item to use its properties
  \item to change its properties (e.g., for animation)
  \item for re-use of common elements (e.g., gradients, images)
  \end{itemize}
\item Used to \textit{create relationships} between elements
\end{itemize}

%%% Show Qt Creator id highlighting feature

\end{slide}

%----------------------------------------------------------------------

\begin{slide}{2009}\frametitle{Using Identities}

\flushedImageDoubleWidth{qml-intro/images/identity.png}
% declarative-uis/concepts/identity.qml
\inputqml{qml-intro/colorized/identity} 

\demo{qml-intro/ex-concepts/identity.qml}

\end{slide}

%----------------------------------------------------------------------

\begin{slide}{2008}\frametitle{Using Identities}

\flushedImageDoubleWidth{qml-intro/images/identity.png}
% declarative-uis/concepts/identity.qml
\begin{qml}
\qtt{\qc{class}{Text}~\{}\\
\qtt{~~~~\qc{type}{id}:~textElement}\\
\qtt{~~~~\qc{lightgray}{x:~50;~y:~25}}\\
\qtt{~~~~\qc{lightgray}{text:~"Qt~Quick"}}\\
\qtt{~~~~\qc{lightgray}{font.family:~"Helvetica";~font.pixelSize:~50}}\\
\qtt{\}}\\
\vspace*{0.5em}
\qtt{\qc{class}{Rectangle}~\{}\\
\qtt{~~~~\qc{lightgray}{x:~50;~y:~75;~height:~5}}\\
\qtt{~~~~\qc{type}{width}:~textElement.width}\\
\qtt{~~~~\qc{lightgray}{color:~"green"}}\\
\qtt{\}}
\end{qml}      

\vspace*{0.5em}
\begin{itemize}
\item \qic{class}{Text} element has the identity, \qtt{textElement}
\item \qic{type}{width} of \qic{class}{Rectangle} bound to \qic{type}{width} of
      \qtt{textElement}
\item Try using \qic{class}{TextInput} instead of \qic{class}{Text}
\end{itemize}

%%% Can identities be changed/reassigned?
%%% Restrictions on IDs? Can we use "parent"?

\end{slide}

%----------------------------------------------------------------------
\begin{slide}{2007}\frametitle{Methods}

\begin{itemize}
\item Most features are accessed via properties
\item Some actions cannot be exposed as properties
\item Elements have methods to perform actions:
  \begin{itemize}
  \item \qic{class}{TextInput} has a \qic{type}{selectAll()} method
  \item \qic{class}{Timer} has \qic{type}{start()}, \qic{type}{stop()}
        and \qic{type}{restart()} methods
  \item \qic{class}{Particles} has a \qic{type}{burst()} method
  \end{itemize}
\item All methods are public in QML
\item Other methods are used to convert values between types:
  \begin{itemize}
  \item \qic{class}{Qt}.\qic{type}{formatDateTime(datetime, format)}
  \item \qic{class}{Qt}.\qic{type}{md5(data)}
  \item \qic{class}{Qt}.\qic{type}{tint(baseColor, tintColor)}
  \end{itemize}
\end{itemize}

\end{slide}

%----------------------------------------------------------------------

\begin{slide}{2006}\frametitle{Basic Types}
Property values can have different types:
\begin{itemize}
\item Numbers (int and real): \qic{number}{400} and \qic{number}{1.5}
\item Boolean values: \qic{number}{true} and \qic{number}{false}
\item Strings: \qic{string}{"Hello Qt"}
\item Constants: \qic{class}{AlignLeft}
\vspace*{0.25em}
\item Lists: \qic{black}{[ ... ]}
  \begin{itemize}
  \item lists with one item can be written as just the item itself
  \end{itemize}
  \vspace*{0.25em}
\item Scripts:
  \begin{itemize}
  \item included directly in property definitions
%  \item can be written inside \qic{black}{\{ ... \}} for clarity
  \end{itemize}
  \vspace*{0.25em}
\item Other types:
  \begin{itemize}
  \item colors, dates, times, rects, points, sizes, 3D vectors, ...
  \item usually created using constructors
  \end{itemize}
\end{itemize}

%%% None/null/nil type?

\doc{qdeclarativeglobalobject.html\#types}{QML Types}
\end{slide}

%----------------------------------------------------------------------
\begin{slide}{2005}\frametitle{Summary}

\begin{itemize}
\item QML defines user interfaces using elements and properties
  \begin{itemize}
  \item elements are the structures in QML source code
  \item items are visible elements
  \end{itemize}
\item Standard elements contain properties and methods
  \begin{itemize}
  \item properties can be changed from their default values
  \item property values can be expressions
  \item \qic{type}{id} properties give identities to elements
  \end{itemize}
\item Properties are bound together
  \begin{itemize}
  \item when a property changes, the properties that reference it are updated
  \end{itemize}
\item Some standard elements define methods
\item A range of built-in types is provided
\end{itemize}

%%% Show property, method syntax

\end{slide}

%%% Scopes - lexical vs. dynamic
                                   

%----------------------------------------------------------------------
\begin{slide}{2004}\frametitle{Exercise}

\begin{enumerate}
\item How do you request features from a certain version of Qt?
% Import Qt with the relevant version.
\item What is the difference between \qic{class}{Rectangle} and \qic{type}{width}?
% Rectangle is an element; width is a property.
\item How would you create an element with an identity?
% Create an id property with a suitable name
\item What syntax do you use to refer to a property of another element?
% identity.property
\end{enumerate}

\end{slide}

%----------------------------------------------------------------------
\begin{slide}{1664}\frametitle{Exercise~\textendash~Items}

% Something about nesting items and relative positions

\flushedImage{qml-intro/images/nested-elements.png}
The image on the right shows two items and two child items inside
a 400 $\times$ 400 rectangle.\\

\begin{enumerate}
\item Recreate the scene using \qic{class}{Rectangle} items.
\item Can items overlap?\\
{\small Experiment by moving the light blue or green\\
rectangles.}
\item Can child items be displayed outside their parents?\\
{\small Experiment by giving one of the child items negative coordinates.}
\end{enumerate}

\end{slide}


\ifthenelse{\boolean{cc-license}}
{
  
  \begin{slide}{999999}
    \vspace{5em}
    \textcopyright~2010 Nokia Corporation and its Subsidiary(-ies).    
    \vspace{2em}                                                     
    \begin{itemize}
      \item[] \footnotesize{The enclosed Qt Training Materials are provided under the
        Creative Commons Attribution-ShareAlike 2.5 License Agreement.}
      \item[]
      \item[] \mbox{\includegraphics[width=20mm]{images/cc-by-sa}}
      \item[]
      \item[] \footnotesize{The full license text is available here:  
      \url{http://creativecommons.org/licenses/by-sa/2.5/legalcode}}
      \item[]    
      \item[] \footnotesize{Nokia, Qt and the Nokia and Qt logos are the registered trademarks 
      of Nokia Corporation in Finland and other countries worldwide.}
    \end{itemize}
    
  \end{slide}
  
} % end if part
{ % else part
  
  \begin{slide}{999999}
    \vspace{5em}
    \textcopyright~2010 Nokia Corporation and its subsidiary(-ies).
    
    \vspace{2em}
    {\small Nokia, the Nokia logo, Qt, and the Qt logo are trademarks of Nokia
      Corporation and/or its subsidiaries in Finland and other countries.
      All other company, product, or service names may be trademarks or service marks
      of others and are the property of their respective owners. The use of the word
      partner does not imply a partnership relationship between Nokia and any other
      company.}
    
  \end{slide}
}


\end{document}
